\section*{About}

{
    Dr. Ariel Noyman is an urban scientist, architect, and urban designer whose work explores how design and technology can \textbf{democratize data-driven urban processes}. Trained and practiced as an architect and urban designer, he bridges urban analytics, simulation, and participatory design across research, practice, and teaching. At MIT, Noyman developed CityScope, an urban human-computer interaction platform for real-time simulation and collaborative planning, deployed in more than twenty cities worldwide. He also helped establish a global network of City Science living laboratories active in over twelve cities. Noyman contributed to New Meadowlands, one of the 2014 Rebuild by Design winners, and to Hamburg's Finding Places project, which URBACT named a Good Practice in 2017. He was on the Chyutin Architects team that won the 2010 Museum of Tolerance design competition. He has published widely and taught urban science, urban design, and architecture design studios at MIT, Cornell, Northeastern University, CUNY, and Bezalel Academy of Arts and Design. He holds a Ph.D. and SM from MIT and a B.Arch., cum laude, from Bezalel Academy.
}

\subsection*{Research Areas}
\textit{Democratizing Urban Processes:} Data-driven approaches to understanding, designing, and collaborative spatial planning.
\\
\textit{Urban Modeling and Simulation:} Developing urban simulation methods focusing on behavior and experience.
\\
\textit{City Science Living-Labs:} Developing a global network of urban research labs.
